\section{The interface}
\label{sec:interface}

\begin{cppcode}
namespace driver::can
{
struct Frame;

class Interface
{
public:
    virtual ~Interface() noexcept = default;

    [[nodiscard]] virtual bool send(const Frame& frame)      = 0;
    [[nodiscard]] virtual bool receive(Frame& frame)         = 0;
    [[nodiscard]] virtual bool hasError() const noexcept     = 0;
    virtual void clearError() noexcept                       = 0;
};
} // namespace driver::can
\end{cppcode}

Every part of it, one at a time:

\textbf{A virtual destructor.} Objects will be deleted through an \code{Interface*}. Without
\code{virtual}, only the base class's destructor runs, and a \code{Spi} holding resources leaks
them --- undefined behaviour, and one of the few bugs in this design that no test catches.
\code{= default} because there is nothing to do; \code{noexcept} because destructors must never
throw.

\textbf{\code{= 0} on all four.} The class is purely abstract: it cannot be instantiated and has
no implementation to inherit by accident. It is a contract and nothing else.

\textbf{\code{[[nodiscard]]} on three of the four.} \code{send()} returns \code{false} when the
hardware is busy. A caller that ignores that believes the frame went out --- the most common
driver bug there is, turned into a compiler warning here.

\textbf{\code{const} on \code{hasError()}.} Asking about an error does not change the driver, and
the method can therefore be called on a \code{const Interface&}.

\textbf{\code{noexcept} on the two error methods.} They are called from error-handling paths, and
an error-handling path that can itself throw is hard to reason about. They also do nothing that
\emph{can} fail.

\subsection{What is not here}

No \code{init()}. No \code{setBaudRate()}. No \code{readRegister()}.

\textbf{The test is: can both \code{Stub} and \code{Spi} implement the method meaningfully?}
\code{readRegister()} they cannot --- the stub has no registers. Such a method would have forced
the stub to lie, and forced the application to know that registers exist.

\section{The stub}
\label{sec:stub}

\code{Stub} is \code{final}, inherits \code{Interface}, and simulates CAN in memory:

\begin{booktable}{@{}lL@{}}
\thead{Method} & \thead{Contract} \\ \midrule
The constructor & \code{myLastFrame} empty, \code{myInitialized = true},
  \code{myHasData = false}. \\
\code{send()} & \code{false} if uninitialised. Otherwise: store the frame, set \code{myHasData},
  return \code{true}. \\
\code{receive()} & \code{false} if uninitialised or without data. Otherwise: copy out the frame,
  clear \code{myHasData}, return \code{true}. \\
\code{hasError()} & Always \code{false}. \\
\code{clearError()} & Does nothing. \\
\code{inject()} & Stores a frame for a later \code{receive()}. \\
\end{booktable}

The copy and move operations are deleted in the public part. A \code{Stub} represents a
(simulated) hardware resource, and two copies of the same resource are not two resources.

\subsection{Why \texorpdfstring{\code{inject()}}{inject()} is not in the interface}

It exists only for the tests' sake: it simulates a frame having arrived on the bus. A \code{Spi}
cannot implement it --- it cannot conjure up traffic that does not exist.

The consequence is the pattern that recurs throughout the book:

\begin{cppcode}
driver::can::Stub stub{};            // concrete type: has inject()
app::EchoNode node{stub, 0x200U};    // takes an Interface&: does not
stub.inject(frame);                  // the test talks to the stub directly
node.run();                          // the component knows nothing about it
\end{cppcode}

The test holds the \textbf{concrete} type, the component sees only the \textbf{interface}.

\subsection{What the stub does not simulate}

No errors. No busy hardware. No queue --- a second \code{send()} overwrites the first. No
arbitration, no CRC.

That is deliberate. The stub's job is to let \code{EchoNode} run, not to be a CAN simulator.

\begin{aside}
\note Knowing where the limit of a test double's realism lies is a skill in itself. \textbf{A stub
that grows until it is a second implementation of the hardware is a stub that needs testing
itself.} Anybody who wants to test how the driver behaves when the hardware is busy needs a
different double, one seam further down --- and that arrives in \kapref{ch:testning}.
\end{aside}

\subsection{Stub, not mock}

A \textbf{stub} answers calls with pre-programmed answers; the test checks what the calling code
did with the answers. A \textbf{mock} also checks that it was called correctly --- the right
method, the right arguments, the right number of times.

\code{Stub} is a stub. \code{ScriptedTransport} in \kapref{ch:testning} is closer to a mock.

\subsection{The seam has already paid off}

The argument for an interface is usually hypothetical. Here it is not: an earlier version of the
design reached the registers through \textbf{memory-mapped I/O} at \code{0xFF200000}, which
assumes an address bus out to the FPGA fabric that the microcontroller does not have.

Everything above \code{Interface} survived the change to SPI \textbf{unchanged}. The old variant
had \code{Fpga} and \code{volatile} pointers where we have \code{Spi} and \code{ByteTransport}.
Neither of them is wrong; they are two implementations under the same seam.
